\section{Introduction}\label{sec:intro}

Scheduling decides which piece of work is executed when, and it is one of the oldest problems in operations research and computer engineering \cite{smith1956,jackson1955,graham1979,pinedo2016}. Classical dispatching rules such as First-In-First-Out (FIFO), Shortest Processing Time (SPT) and Earliest Deadline First (EDF) are attractive because they are transparent, cheap to compute and, under the right assumptions, provably optimal: EDF minimises maximum lateness on a single resource with constant capacity \cite{jackson1955,liu1973}, and SPT minimises mean flow time \cite{smith1956}. The common assumption underlying these guarantees is that the resource executing the work has a \emph{fixed, state-independent processing capacity}. This assumption is reasonable for processors and machines but is violated whenever the resource is a person.

Human work capacity is dynamic. Sustained effort depletes attentional resources and produces the vigilance decrement first documented by Mackworth~\cite{mackworth1948} and since replicated across many tasks \cite{warm2008,thomson2015,esterman2019}; mental fatigue accumulates with time-on-task and degrades both speed and accuracy \cite{boksem2008,lorist2000,hockey2013,ackerman2011}; and alertness follows a circadian rhythm with a pronounced early-afternoon dip \cite{akerstedt1997,monk2005,schmidt2007,blatter2007}. Rest breaks partially reverse these effects, and field studies consistently report that well-timed breaks preserve or even increase productivity despite the time they consume \cite{henning1997,dababneh2001,galinsky2000,tucker2003,albulescu2022}. A schedule that is optimal for a constant-capacity machine can therefore be poor for a human: it front-loads demanding work regardless of cognitive state, never allocates recovery, and drives the worker into a low-capacity regime in which deadlines are missed precisely because capacity was ignored. The concern is not only productivity. Chronic fatigue is associated with errors and accidents \cite{dawson1997,williamson2011}, and the move from Industry~4.0 towards a human-centric Industry~5.0 explicitly calls for production and service systems that adapt to the well-being of their workers rather than the reverse \cite{breque2021,xu2021,lu2022,leng2022,ivanov2023}.

Operations-management research has long recognised this gap. Models of worker fatigue and recovery have been integrated into dual-resource-constrained systems \cite{jaber2010,jaber2013,givi2015}, rest-allowance and break-scheduling problems have been formulated analytically \cite{bechtold1984,konz1998,calzavara2019}, and taxonomies for merging scheduling theory with human factors have been proposed \cite{lodree2009,neumann2010,boudreau2003,katiraee2021}. These approaches, however, usually assume that the worker's state is known or deterministic and optimise an offline plan; they do not provide an adaptive \emph{policy} that reacts to the evolving, uncertain state of a specific person during the day. Deep reinforcement learning (DRL), in turn, has produced effective adaptive schedulers for job shops, clusters and production lines \cite{mao2016,mao2019,waschneck2018,luo2020,zhang2020l2d,park2021,song2023}, and recent surveys document rapid growth of the field \cite{kayhan2023,panzer2022,khadivi2025}. Yet the resources in these studies are machines, and the few recent studies that include worker fatigue embed it in large shop-floor models where fatigue is a deterministic, fully observed variable \cite{tan2021,hu2025}. Two features of the human problem are thus rarely addressed together: \emph{the decision of when to rest is itself a scheduling action}, and \emph{the cognitive state that should drive that decision cannot be measured directly}. In practice, attention and fatigue are inferred from noisy proxies such as physiological signals, eye tracking, performance probes or self-report \cite{charles2019,borghini2014,hogervorst2014,basner2011,hart1988}, which makes the scheduling problem partially observable.

This paper studies attention-aware task scheduling for a single knowledge or production worker as a partially observable Markov decision process (POMDP) \cite{astrom1965,smallwood1973,kaelbling1998} and solves it with a memory-augmented deep Q-network. We pay particular attention to the methodological issues that make DRL comparisons hard to trust \cite{henderson2018,agarwal2021}: we compare against strong, \emph{tuned}, cognition-aware heuristics rather than only against FIFO and EDF; we use disjoint seed sets for training, tuning and testing; we evaluate on 500 held-out workdays with common random numbers; we report ten independent training seeds with confidence intervals and multiplicity-corrected paired tests; and we separate effects of the reward function from effects of the dynamics through ablation.

The contributions are as follows.
\begin{enumerate}[label=(\roman*)]
  \item \textbf{A cognitively grounded scheduling simulator.} We formulate a workday-scale scheduling environment in which a latent homeostatic attention resource, an exponentially accumulating and recovering fatigue state \cite{jaber2010}, and a circadian process \cite{akerstedt1997,monk2005} jointly determine the worker's instantaneous work rate. Tasks arrive stochastically with heterogeneous difficulty, work content and deadlines. Rest is an explicit action. The simulator, its parameters and their behavioural rationale are fully specified and released.
  \item \textbf{A POMDP formulation with a multi-objective reward.} The scheduler observes only noisy estimates of attention and fatigue together with a compact, size-invariant summary of the task queue. The reward combines deadline-weighted task value with explicit attention and fatigue terms, and we analyse the resulting productivity--fatigue trade-off as a Pareto front \cite{roijers2013,hayes2022}.
  \item \textbf{An attention-aware DQN (AA-DQN).} We combine a finite observation history, which acts as an approximate sufficient statistic for the belief state \cite{singh1994,hausknecht2015,ni2022}, with Double Q-learning \cite{vanhasselt2016} and a dueling architecture \cite{wang2016dueling}. A component ablation quantifies the contribution of each element.
  \item \textbf{A rigorous empirical study.} Against seven baselines, including grid-tuned threshold and attention-matching heuristics that observe the same noisy cognitive signals, we test (a) whether learning helps, (b) \emph{why} it helps, through environment ablations that remove observation noise, the circadian process, the cognitive dynamics or the cognitive reward terms, (c) what behaviour is learned, and (d) how the learned policy generalises to unseen worker profiles, workloads and sensor-noise levels.
\end{enumerate}

The remainder of the paper is organised as follows. Section~\ref{sec:related} reviews related work and positions the study. Section~\ref{sec:problem} formulates the problem. Section~\ref{sec:method} presents AA-DQN. Section~\ref{sec:setup} details the experimental protocol. Section~\ref{sec:results} reports results, Section~\ref{sec:discussion} discusses implications, ethics and limitations, and Section~\ref{sec:conclusion} concludes.
